\begin{table}[tbp]
\centering
\caption{补充：潜状态替代方案比较}
\label{tab:s06_q1_latent}
\footnotesize
\setlength{\tabcolsep}{3.5pt}
\begin{tabular}{llrrrr}
\toprule
潜状态方法 & 划分 & Δ vs PCA & 95\% CI 下 & 95\% CI 上 & 正增量维数 \\
\midrule
pca & 随机 CV & 0.0000 & 0.0000 & 0.0000 & 0 \\
pca & 分组 CV & 0.0000 & 0.0000 & 0.0000 & 0 \\
fa & 随机 CV & -0.0149 & -0.0308 & 0.0006 & 22 \\
fa & 分组 CV & -0.0055 & -0.0194 & 0.0080 & 25 \\
pls8 & 随机 CV & 0.0117 & 0.0026 & 0.0213 & 33 \\
pls8 & 分组 CV & 0.0131 & 0.0042 & 0.0215 & 34 \\
pls12 & 随机 CV & -0.0015 & -0.0113 & 0.0091 & 28 \\
pls12 & 分组 CV & -0.0115 & -0.0213 & -0.0008 & 20 \\
cca8 & 随机 CV & -2.0636 & -2.4601 & -1.7210 & 0 \\
cca8 & 分组 CV & -2.5898 & -2.9357 & -2.2627 & 0 \\
\bottomrule
\end{tabular}
\end{table}
% 注（写入图注/表注，勿直接编译）：以 14 维化学 PCA 为参照，比较因子分析（FA）与偏最小二乘（PLS）潜状态的增量；正增量维数指 50 维中 Δ > 0 的维数。
